\documentclass[12 pt]{exam}
\usepackage[margin=0.75in]{geometry}
\usepackage{amsmath,amssymb,color,amsthm}
\usepackage{enumerate,graphicx,multicol}
\usepackage{wrapfig}
\usepackage{pgf,tikz, subcaption}
\usetikzlibrary{arrows}

\newtheorem{claim}{Claim}
\newtheorem{lemma}{Lemma}

\printanswers

\pagestyle{head}                       % This causes the exam to only have headers, no footers.
%\runningheadrule                     % This causes the headings to have an underline.

%\firstpageheader{Name:\enspace\hbox to 4.5in{\hrulefill} Section:\enspace\hbox to 1.2in{\hrulefill}}{}{}
%\extraheadheight{.25in}

\begin{document}
\hrule
\begin{center}
  \huge{Math 2710 Problem Set 7: Ch 7}    %Title
\end{center}
\hrule
\vskip .2in


\begin{questions}
\question Suppose $a \in \mathbb{Z}$. Prove that $14 \mid a$ if and only if $2 \mid a$ and $7 \mid a$.

\begin{solution}\begin{proof}
  If $14 | a$ then $2 | a$ and $7 | a$.

  Let $a \in \mathbb{Z}$. Assume that $14 | a$. Then $a = 14n$ where $n \in \mathbb{Z}$.
  Then $a = 2(7n)$. Since $7n \in \mathbb{Z}$, then $2 | a$. Also, $a = 7(2n)$, and
  since $2n \in \mathbb{Z}$, then $7 | a$. So if $14 | a$ then $2 | a$ and $7 | a$.

  Let $a \in \mathbb{Z}$. Assume that $2 | a$ and $7 | a$. Then $a = 2n$ and $a = 7k$ where $n, k \in \mathbb{Z}$.
  Then $2n = 7k$. Since $2n$ is even, then $7k$ must be even, so $k$ must be even. Then $k = 2m$, where
  $m \in \mathbb{Z}$. Then $a = 7(2m) = 14m$, where $m \in \mathbb{Z}$, so $14 | a$. So if $2 | a$ and $7 | a$, then $14 | a$. 

  Since we have show that $(14 | a) \Rightarrow (2 | a) \land (7 | a)$ and $(2 | a) \land (7 | a) \Rightarrow
  (14 | a)$, then $(14 | a) \iff (2 | a) \land (7 | a)$.
\end{proof}\end{solution}


\question 
\begin{parts}
\item Prove that there exists a positive real number $x$ for which $$x^2 < \sqrt{x}.$$

\begin{solution}\begin{proof}
  Let $x = \frac{1}{4}$. Then $x^2 = \frac{1}{16}$ and $\sqrt{x} = \frac{1}{2}$, so $x^2 < \sqrt{x}$ since
  $\frac{1}{16} < \frac{1}{2}$. Since $\frac{1}{4}$ is a positive real number that satisfies $x^2 < \sqrt{x}$,
  then the statement is true.
\end{proof}\end{solution}

\item For what statement is the example you gave above a counterexample?

\begin{solution}
  $\forall x \in \mathbb{R}, x^2 \ge \sqrt{x}$.
\end{solution}


\end{parts}

\question Give a counterexample to disprove each of the following false quantified statements:
\begin{parts}
\part For any sets $A$ and $B$, $\mathcal{P}(A)-\mathcal{P}(B)\subseteq\mathcal{P}(A-B)$.

\begin{solution}\begin{proof}
  Let $A = \{1,2\}$ and $B = \{2\}$. Then $\mathcal{P}(A) = \{\emptyset, \{1\}, \{2\}, \{1,2\}\}$
  and $\mathcal{P}(B) = \{\emptyset, \{2\}\}$. So $\mathcal{P}(A) - \mathcal{P}(B) =  \{\{1\}, \{1,2\}\}$.
  Also, $\mathcal{P}(A-B) = \mathcal{P}(\{1\}) = \{\emptyset, \{1\}\}$.
  Since $\{1, 2\} \in \mathcal{P}(A) - \mathcal{P}(B)$ and $\{1, 2\} \notin \mathcal{P}(A-B)$, then
  $\mathcal{P}(A) - \mathcal{P}(B) \subsetneq \mathcal{P}(A-B)$ and the statement is fase.
\end{proof}\end{solution}


\part For every positive real number $a$, there is a positive real number $b$ such that $b<a$ and if $c\in \mathbb{R}^+$, then $bc\geq c$.

\begin{solution}\begin{proof}
Let $a = 1$. For the sake of contradiction assume that there exists a positive real $b$ such that $b < 1$
and if $c \in \mathbb{R}^+$, then $bc \ge c$. If $b < 1$, then $bc < c$. But then we have that both
$bc \ge c$ and $bc < c$. This is a contradiction so the statment is false.
\end{proof}\end{solution}


\end{parts}

\question Prove each of the following existence statements:
\begin{parts}
\part For every non-zero real number $x$, there exists a real number $y$ such that $x+y=7$.

\begin{solution}\begin{proof}
Let $x$ be a non-zero real number. Choose $y = 7 - x$. Then $x + y = x + (7 - x) = 7$.
\end{proof}\end{solution}

\part There exists a natural number $M$ such that for all natural numbers $n>M$, $\dfrac{1}{n}<0.105$.

\begin{solution}\begin{proof}
  Let $M = 10$. Let $n \in \mathbb{N}$ and $n > M$, so $n > 10$. Then $\frac{1}{n} < \frac{1}{10}$.
  Since $\frac{1}{10} < 0.105$, then $\frac{1}{n} < 0.105$.
\end{proof}\end{solution}

\end{parts}


\question Disprove each of the following existence statements:
\begin{parts}
\part There exists a largest natural number $M$.

\begin{solution}\begin{proof}
Assume for the sake of contradiction that there exists a largest natural number $M$. Consider the number
$M + 1$. Since $M+1 > M$ and $M + 1 \in \mathbb{N}$, we have found a contradiction. We assumed that $M$ was
the largest natural number but found that $M+1$ is a larger natural number. So there is no largest natural number.
\end{proof}\end{solution}

\part There exists a rational number $x$ and an irrational number $y$ such that $x+y$ is rational (Hint: first prove that the sum of two rational numbers is rational).

\begin{solution}\begin{proof}
  Let $x \in \mathbb{Q}$ and $y \in \mathbb{R-Q}$. Assume for the sake of contradiction that $x+y \in \mathbb{Q}$.
  Then $x = \frac{p}{q}$ where $p,q \in \mathbb{Z}$ and $q \ne 0$, and $x+y = \frac{m}{n}$, where $m,n
  \in \mathbb{Z}$ and $n \ne 0$. Then $y = \frac{m}{n} - \frac{p}{q} = \frac{mq - pn}{nq}$. Since $mq-pn, nq \in \mathbb{Z}$ and $nq \ne 0$ since $n \ne 0$ and $q \ne 0$, then $y \in \mathbb{Q}$. This is a contradiction
  since was assumed that $y \in \mathbb{R-Q}$ but we found that $y \in \mathbb{Q}$. Therefore there does not
  exist a rational $x$ and irrational $y$ such that $x+y$ is rational.
\end{proof}\end{solution}
\end{parts}





\section*{Challenge Problems}
\question (If you have taken Math 2210Q) Prove at least four parts of the Invertible Matrix Theorem are equivalent.

Prove each of the following statements.

\question Suppose $a, b \in \mathbb{N}$. Then $a=\gcd(a,b)$ if and only if $a \mid b.$
\begin{solution}\begin{proof}
  Let $a, b \in \mathbb{N}$.

  Assume that $a = \gcd(a,b)$. Then $a | b$. So if $a = \gcd(a,b)$, then $a | b$.

  Assume that $a | b$. Assume for the sake of contradiction that $a \ne \gcd(a, b)$. Then there must exist some
  $k \in \mathbb{N}$ such that $k = \gcd(a, b)$ and $k > a$. Since $k = \gcd(a, b)$, then $k | a$.
  This is a contradiction since no natural is divisible by a
  natural greater than itself, so we have that $a | b \Rightarrow a = \gcd(a,b)$.

  Since we have shown that $a = \gcd(a, b) \Rightarrow a | b$ and $a | b \Rightarrow a = \gcd(a,b)$,
  then $a = \gcd(a,b) \iff a | b$.
\end{proof}\end{solution}

\question Suppose $a, b \in \mathbb{N}$. Then $a=lcm(a,b)$ if and only if $b \mid a.$
\begin{solution}\begin{proof}
  Let $a, b \in \mathbb{N}$

  Assume that $a = lcm(a, b)$. Then $b | a$. So we have that $a = lcm(a, b) \Rightarrow b | a$.

  Assume that $b | a$. Assume for the sake of contradiction that $a \ne lcm(a,b)$. Then there must exist some
  $k \in \mathbb{N}$ such that $k = lcm(a, b)$ and $k < a$. Since $k = lcm(a, b)$, then $a | k$.
  But this is a contradiction since no natural divides another
  natural lower than itself. So we have that $b | a \Rightarrow a = lcm(a, b)$.

  Since we have shown that $a = lcm(a, b) \Rightarrow b | a$ and $b | a \Rightarrow a = lcm(a,b)$,
  then $a = lcm(a,b) \iff b | a$.
\end{proof}\end{solution}


\question If $p >1$ is an integer and $n \nmid p$ for each integer $n$ for which $2 \leq n \leq \sqrt{p}$ then $p$ is prime. 

\end{questions}
\end{document}
